\section{An exact complex edge sum}\label{cta:edge}

\subsection{The interpolation and its normalization}
We use the canonical factorial lemma
\cite[Lemma~2.1]{ConwayAnalysis}, whose statement we record.
On \(-1<\Rea z<N\),
\begin{equation}\label{cta:G}
 G_N(z)=\left(\frac{2\pi}{N}\right)^z\Gamma(z+1)
       \prod_{a\ge1}
       \frac{\Gamma(Na+z+1)}
            {\Gamma(Na-z)(Na)^{2z+1}}
\end{equation}
is holomorphic and nonzero, with positive-base powers.
It satisfies \(G_N(0)=1\), \(G_N(k)=S_k\), and
\[
 \frac{G_N(z+1)}{G_N(z)}=2\sin\frac{\pi(z+1)}N,\qquad
 G_N(z)G_N(d-z)=N,\qquad G_N(s)G_N(-s)=H_N(s)^{-1}.
\]
On every compact subset of \(0<\Rea x<1\),
\begin{equation}\label{cta:G-asymptotic}
 \log G_N(Nx)=NJ(x)+\tfrac12\Log(2N\sin\pi x)
             +\frac{\pi\cot\pi x}{12N}+O(N^{-3}).
\end{equation}
All fixed derivatives, and any fixed further inverse-power order,
are uniform on smaller compact sets. Logarithms and square roots
are continued from the positive real interval. This numbered
lemma includes the convergent product and summable Stirling proof;
no periodic factor may be inserted into \eqref{cta:G}.
For background on those classical expansions see \cite[Section~5.11]{DLMF}.

The shift equation now proves the exact interpolation needed here.
For \(u=Nt/(2\pi)\), define
\begin{equation}\label{cta:interpolant}
 \mathcal B_N(x,t)=
 \frac{G_N(d/2+u)G_N(d/2-u)}
 {G_N(d/2-Nx+u)G_N(d/2-Nx-u)G_N(Nx)^2}.
\end{equation}
At \(x=k/N\) this equals \(e_k(-2\cos t)\).
Indeed successive quotients are the original cosine numerator
and sine-factorial denominator of \eqref{cta:basis}.
The exact edge insertion is
\begin{equation}\label{cta:insertion}
 \mathcal I_N(x,s)=
 \frac{G_N(Nx)^2}{G_N(Nx-s)G_N(Nx+s)},\qquad
 \mathcal I_N(k/N,s)=\frac{H_N(s)}{W_k(s)}.
\end{equation}
The second identity follows by iterating the same shift equation
and using \(G_N(s)G_N(-s)=H_N^{-1}\).

Put
\begin{equation}\label{cta:b}
 b(x,t)=\int_0^x\Log(2\cos(2\pi v)+2\cos t)\,\dd v-2J(x).
\end{equation}
Branches start at real \(x\in I,t=0\).
On a fixed complex tube about \(I\), \eqref{cta:G-asymptotic} gives
\begin{equation}\label{cta:bulk}
 \begin{split}
 \mathcal B_N(x,t)
   &=\frac{e^{Nb(x,t)}}{2N\sin\pi x}
          (1+N^{-1}b_1(x,t)+O(N^{-2})),\\
 \mathcal I_N(x,s)
   &=1-\frac{\pi s^2}{N}\cot\pi x+O_R(N^{-2}).
 \end{split}
\end{equation}
The first correction \(b_1\) is independent of \(s\).
To check the leading constant, each argument \(d/2\pm u\)
has the shift \(-1/2\), so its square-root factor is \(\sqrt N\);
the analogous denominator factors cancel these, leaving
\((2N\sin\pi x)^{-1}\). Thus no midpoint shift has been discarded.
For the second identity use
\((\log G_N)''(Nx)=\pi\cot(\pi x)/N+O(N^{-2})\)
in the symmetric difference. All remainders are holomorphic
with uniform fixed derivatives on a smaller tube.

\subsection{The exact stationary phase and an available contour}
For \(A(x,t,u)=b(x,t)+b(x,u)\), its critical branch satisfies
\begin{equation}\label{cta:critical-wire}
 \cos(2\pi x)=\frac{1-\cos t\cos u}{2+\cos t+\cos u},
 \qquad x(0,0)=1/4.
\end{equation}
This is equivalent to \eqref{cta:wire}. It follows by writing
\(A_x=0\), squaring the products of the nonzero factors, and
choosing the branch at zero; the original unsquared equation
holds there and hence nearby.
Substitute \(z=\tan\pi v\) into the \(t\)-derivative of
\eqref{cta:b}. At its critical \(x\), the resulting logarithmic
antiderivative is
\(-\pi^{-1}\operatorname{arsinh}(\tan(t/2)\cos(u/2))\).
Thus \(\psi=\stat_x A\) is exactly \eqref{cta:psi-definition}.
In particular \(\psi\) is separately even.

We need a domain for the finite sum, not just this limiting phase.
On \(|t|,|u|\le1/4\), \(|\tan(t/2)|<0.128\); the tangent form
of \eqref{cta:critical-wire} implies
\[
 |\cos(2\pi x)|<0.017,\qquad |x(t,u)-1/4|<0.003.
\]
On \(\Rea x\in I,\ |\Ima x|\le1/50\), direct differentiation gives
\begin{equation}\label{cta:Axx}
 A_{xx}=\sum_{\xi=t,u}
 \left[-\frac{2\pi\sin2\pi x}{\cos2\pi x+\cos\xi}
                         -2\pi\cot\pi x\right],
 \qquad |A_{xx}|<44,\quad \Rea A_{xx}<-12.
\end{equation}
Here are sufficient elementary bounds: \(|\cos2\pi x|<0.338\),
\(|\cos\xi-1|<0.032\), \(|\sin2\pi x|<1.009\), and
\(\Rea\cot\pi x>0.72\).
The real part of
\(\sin2\pi x/(\cos2\pi x+\cos\xi)\) exceeds \(0.328\):
its conjugated numerator is at least
\(0.95(0.655)-0.040(0.159)\), and its squared denominator is
at most \(1.37^2\). These also prove the modulus bound.

Enclose precisely the retained integers \(k\) by a rectangle in
the \(x\)-plane with vertical half-grid sides and horizontal
height \(h=1/500\). The residue formula with
\(\pi N\cot(\pi Nx)\) equals the finite sum exactly.
On a vertical half-grid side the cotangent has modulus at most
one. On horizontal sides it equals \(\mp\ii+O(e^{-2\pi Nh})\).
If \(X=\Rea(x-x_*),Y=\Ima(x-x_*)\), Taylor's formula and
\eqref{cta:Axx} imply
\begin{equation}\label{cta:contour-bound}
 \Rea\{A(x)-A(x_*)\}\le-6X^2+44|XY|+22Y^2.
\end{equation}
On a vertical side, for all large \(N\),
\(|X|\ge0.047-o(1)\) and \(|Y|\le0.005\);
the right side is below \(-0.002+o(1)\).
On a horizontal side its maximum in \(X\) is at most
\((308/3)Y^2\le77/30000\), less than
\(2\pi h=\pi/250\). All nonconstant cotangent modes and
vertical currents therefore have a fixed exponential loss
relative to \(e^{N\psi}\).

The constant horizontal terms give \(N\) times a real integral.
Translate this integral to the real tangent through \(x_*\).
Its endpoint connectors have the same loss, and
\(\Rea A_{xx}<-12\) gives a nondegenerate Gaussian. The
ordinary Taylor proof of the saddle expansion applies uniformly:
outside a fixed small saddle interval there is a fixed loss;
inside it put \(x-x_*=N^{-1/2}z\), expand, and integrate
the Gaussian-dominated remainder. Odd terms integrate to zero.
This also gives any fixed order and its parameter derivatives.
All Gamma arguments in \eqref{cta:interpolant} stay in regular
bulk strips throughout these contours.

It follows that
\begin{equation}\label{cta:edge-expansion}
 \begin{split}
 K_N(t,u;0)&=N^{-3/2}a_e(t,u)e^{N\psi(t,u)}
                     (1+N^{-1}a_{e,1}(t,u)+O(N^{-2})),\\
 a_e&=\frac{\sqrt{2\pi}}
       {4\sin^2(\pi x_*)\sqrt{-A_{xx}(x_*,t,u)}} .
 \end{split}
\end{equation}
The square root is positive at zero and continued on the disk;
no factor vanishes there. In particular \(a_e(0,0)=1/4\).
Applying the same contour to \eqref{cta:insertion} proves
\begin{equation}\label{cta:packet-expansion}
 \begin{split}
 B_N(\boldsymbol\tau,0)&=
   N^{-3e/2}a(\boldsymbol\tau)e^{N\Psi(\boldsymbol\tau)}
           (1+N^{-1}a_1(\boldsymbol\tau)+O(N^{-2})),\\
 \frac{B_N(\boldsymbol\tau,s)}{B_N(\boldsymbol\tau,0)}
   &=1-\frac{\pi s^2}{N}\sum_{vw\in E}
             \cot(\pi x(\tau_v,\tau_w))+O_R(N^{-2}),
 \qquad a=\prod_e a_e .
 \end{split}
\end{equation}
For a holomorphic insertion, its first saddle correction is
\(O(N^{-1})\), since the linear Gaussian term is odd.
Multiplication by its prefactor \(1/N\) in
\eqref{cta:bulk} therefore makes the last remainder \(O(N^{-2})\).
The common \(a_1\) cancels in this ratio.

\section{Angular bounds and weighted curvature}\label{cta:angular}

\subsection{A bound on the entire real angular period}
Let \(N\ge160\), \(k/N\in I\), and
\(r_j=\cos((2j+1)\pi/N)\). The exact product is
\[
 \frac{e_k(-2\cos z)}{e_k(-2)}
   =\prod_{j<k}Q_{r_j}(z),\qquad
 Q_r(z)=\frac{r+\cos z}{1+r}.
\]
Reserve the first \(m_0=\lfloor N/32\rfloor\) factors, where
\(r_j>\cos(\pi/16)>9/10\).
There are at most \(N/20+1\) negative parameters, all \(>-2/5\).
Pair them injectively with parameters among the next
\(N/20+1\) positions, which are \(>7/10\); these positions
are disjoint from the negative ones since \(m_0+N/20+1\le N/8\).
All remaining unpaired parameters are nonnegative.

For a pair \(r>-2/5,s>7/10\), the coefficients of
\[
 (z^2+2rz+1)(z^2+2sz+1)
\]
are \(1,2(r+s),2+4rs,2(r+s),1\), all positive.
Normalized at \(z=1\), with its central Laurent power removed,
it is a symmetric probability measure on
\(\{-2,-1,0,1,2\}\).
A remaining single factor has the analogous interpretation
on \(\{-1,0,1\}\). Its complex-angular modulus is at most
\(\cosh(by)\le e^{b^2y^2/2}\), where \(b\) is its Laurent degree.

For a reserved factor, put \(w=\sin^2(t/2)\). The two linear
pieces of \(|1-2w/(1+r)|\) give
\(|Q_r(t)|\le1-w/2\). For \(|y|\le1/4\),
\[
 |r+\cos(t+\ii y)|^2
 =(r+\cos t)^2+2r\cos t(\cosh y-1)+\sinh^2 y .
\]
Since \(2(\cosh y-1)+\sinh^2y\le4y^2\), each reserved
factor is bounded by \(\exp(-w/4+2y^2)\).
Use \(m_0\ge N/40\), then multiply the paired and reserved
estimates:
\begin{equation}\label{cta:raw}
 \left|\frac{e_k(-2\cos(t+\ii y))}{e_k(-2)}\right|
 \le\exp N\left(-\frac{\sin^2(t/2)}{160}+4y^2\right).
\end{equation}
This pairing precedes absolute values. Individual Laurent
factors need not have nonnegative coefficients.

Set \(a_v=d_v/(2\pi)\). On \(|t_v|\le\pi,\ |y_v|\le1/4\),
\(\sin^2(t_v/2)\ge t_v^2/\pi^2\), so \eqref{cta:raw},
the bounded edge insertions, and the positive zero-angle sum give
\begin{equation}\label{cta:global-Gaussian}
 |B_N(\mathbf t+\ii\mathbf y,s)|
 \le C_RN^A\exp N\left\{
 F_0-\frac1{400}\sum_va_vt_v^2+70\sum_va_vy_v^2\right\}.
\end{equation}
The stronger incidence constants are \(1/(80\pi)>1/400\)
and \(8\pi<70\).
Endpoint translations \(O_R(N^{-1})\) and derivatives cost
only polynomial factors for the fixed tree; one may apply
Cauchy's formula on radius \(1/N\) disks. This is a bound on
the entire angular period, not only near its maximum.

\subsection{The local nonlinear terms}
With
\(Q=\sqrt{1-\sin^2(t/2)\sin^2(u/2)}\), differentiation gives
\begin{equation}\label{cta:edge-Hessian}
 \psi_{tt}=-\frac{\cos(u/2)}{2\pi\cos(t/2)Q},\qquad
 \psi_{tu}=\frac{\sin(t/2)\sin(u/2)}{2\pi Q},\qquad
 \det\psi''=\frac1{4\pi^2}.
\end{equation}
Consequently
\[
 \Psi=F_0-\frac12\sum_v a_v\tau_v^2+\calE(\boldsymbol\tau).
 \]
On the polydisk \(|\tau_v|\le1/4\), put
\(\delta=1/50,\ \varepsilon=1/25\). Then
\begin{equation}\label{cta:error}
 \left|\frac{\calE_v}{\tau_v}\right|\le\delta a_v,\qquad
 \sum_w|\calE_{vw}|\le\varepsilon a_v,\qquad
 |\calE|\le\frac{\varepsilon}{2}\sum_v a_v|\tau_v|^2 .
\end{equation}
The quotient extends at zero by separate evenness.
For explicit bounds write
\(k=\cosh(1/8)-1<0.008\), \(\sigma=\sinh(1/8)<0.126\),
and \(C=(1-\sigma^4)^{-1/2}\).
The normalized diagonal error in \eqref{cta:edge-Hessian}
is at most \(2kC/(1-k)+C-1<1/50\), and the off-diagonal
bound is \(\sigma^2C\); their sum is \(<1/25\).
Summing over incident edges proves the row estimate.
Integrate the diagonal error while fixing the other variables
to prove the first estimate. For the last, sum the edge
quadratic forms and use \(2|zw|\le|z|^2+|w|^2\).
No estimate grows with the number of vertices.
